\section{Limitations and future work}
\label{sec:limitations}

\paragraph{What the graph does not state cannot be preserved.}
The context trees and the coverage objective only reflect the relationships encoded in the ontology.
This matters most for primitive concepts, whose relationships do not fully define them (Section~\ref{sec:snomed}), and SNOMED~CT has many of them: 64\% of the 420,716 concepts of the release we use are primitive, and in the HERO data, 42\% of the 11,796 concepts that belong to the release are primitive and carry 88\% of the events of these 11,796 concepts.
When a primitive concept is not itself a token, its tree loses the part of its meaning that no relationship states; \textit{hyperthyroidism} and \textit{sick-euthyroid syndrome}, two primitive concepts with the same relationships, receive the same tree (Section~\ref{sec:results:examples}).
The status-dependent decay of Eq.~\eqref{eq:coverage-status} addresses this directly, since it charges exactly the hops where meaning can be lost; evaluating it is our next step.
Two further simplifications come from our design rather than from the graph: all relationship types weigh the same in the objective, and context trees do not record how SNOMED~CT groups the attributes of a concept.
Factors that depend on the relationship type fit the same framework (Section~\ref{sec:method:score}).

\paragraph{Transductive and greedy selection.}
The selection optimizes the coverage of a fixed set of concepts.
A new concept can still be tokenized with an existing vocabulary, since Algorithm~\ref{alg:expand} applies to any concept of the graph, but the vocabulary was not chosen to cover it; a new institution or release may require a new run, which is inexpensive and needs no patient data.
The selection also never removes a token, so hubs selected early remain even when later, more specific tokens cover the same concepts.
The guarantee of Eq.~\eqref{eq:guarantee} bounds the resulting loss in coverage, but not in the other metrics; a pass that swaps out tokens with a small marginal loss could reduce this redundancy.

\paragraph{No downstream or external validation.}
We evaluate the tokenization intrinsically, with the metrics of Section~\ref{sec:method:metrics} and with semanticians on 100 concepts, not on a downstream clinical task.
Whether a foundation model trained on our vocabulary predicts clinical outcomes better than one trained on the full or truncated vocabulary remains to be tested by pretraining experiments.
All data also come from a single institution, whose mapping practice and concept distribution may differ elsewhere.
Where mappings between ontologies exist, the same mechanism could also write the concepts of one terminology with tokens of another.

\paragraph{Graph assumptions.}
The method assumes that edges point from a concept to the concepts that describe it, and that the graph is acyclic (Section~\ref{sec:method:problem}).
Because every traversal is bounded by $D$, it still terminates on a graph with cycles, but on a vocabulary with cycles or symmetric relationships, or on a general knowledge graph, the tokens reached from a concept need not describe it.

\section{Conclusion}
\label{sec:conclusion}

We presented an ontology-based tokenizer that represents clinical concepts with a compact vocabulary of ontology concepts, selected by distance-aware greedy coverage and used through context trees.
On 46,150 mapped concepts, it ranked first among 19 configurations at every compression level, and on real EHR data it represented every concept of the ontology graph, including the rare concepts that frequency truncation discards, with a mean proximity-coverage score of 0.88 for these rare concepts.
\needresults{one sentence on the expert result}
Whether these vocabularies improve EHR foundation models is the next test.
